% Source presentations for register row N-x-015.

\subsection{Probability logic and quantitative model checking as separate source presentations}
\label{sec:p09-presentation-n-x-015}

Quantitative reasoning here has three different semantic objects: a
probability assigned to a propositional event, a probability measure on paths
of a discrete transition model, and a probability measure on continuous-time
trajectories.  Definitions~\ref{def:p09-likelihood-logic},
\ref{def:p09-dtmc-mdp}, and~\ref{def:p09-ctmc-csl} give those objects before any
algorithm is stated.

\subsubsection{Halpern's uncertainty chapter and the adjacent linear logic}

Halpern's 2003 first edition places measurable probability, other quantitative
representations of likelihood, relative likelihood, knowledge and probability,
independence, and expectation in separate sections of Chapter~7.  This is a
contents-level architectural fact: a qualified contents scan matching the
first-edition pagination gives the section starts, but no technical body.  The currently visible
Internet Archive scan is instead the access-restricted 2005 paperback, which
Halpern describes as slightly revised; it is not evidence for the 2003 wording.
Accordingly, no grammar, semantic tuple, axiom system, theorem, proof, or
example below is attributed to the unread 2003 chapter.  The technical theorem
is stated only at the strength verified in the adjacent primary paper
\cite{Halpern2003}.

The technical theorem used in
Theorem~\ref{thm:p09-likelihood-logic} instead comes from the adjacent primary
paper of Fagin, Halpern, and Megiddo.  Its terms are integer linear combinations
of expressions \(w(A)\), and formulas are Boolean combinations of linear
inequalities.  Probability structures make every propositionally definable
event measurable.  The axioms combine propositional logic, valid linear
inequality reasoning, normalization, nonnegativity, invariance under
propositional equivalence, and the splitting equation
\[
 w(A\land B)+w(A\land\neg B)=w(A).
\]
The source calls the resulting system
\(\mathsf{AX}_{\mathrm{MEAS}}\).  Its Theorem~2.2 proves soundness and
completeness; Lemma~2.1 derives \(w(\bot)=0\), and disjoint finite additivity
follows from the splitting equation and propositional equivalence
\cite[\S\S2--3]{FaginHalpernMegiddo1990}.

For a concrete calculation, suppose a model satisfies
\[
 w(p)=\frac23,
 \qquad w(q)=\frac12,
 \qquad w(p\land q)=\frac13.
\]
Each equality abbreviates two permitted linear inequalities.  The disjoint
decomposition
\[
 p\lor q=(p\land\neg q)\ \dot\lor\ (p\land q)\ \dot\lor\
           (\neg p\land q)
\]
and finite additivity give
\[
 w(p\lor q)=w(p)+w(q)-w(p\land q)=\frac56.
\]
No independence premise was used.  These particular numbers happen also to
satisfy the independence equation \(w(p\land q)=w(p)w(q)=1/3\), but that
equation played no role.

The completeness proof first chooses a consistent conjunction from a
disjunctive normal form.  If \(p_1,\ldots,p_n\) are the primitive propositions
occurring in it, an \(n\)-atom is a conjunction
\(L_1\land\cdots\land L_n\) in which each \(L_i\) is either \(p_i\) or
\(\neg p_i\).  It is a propositional formula representing one complete truth
assignment, not a new kind of world.  Every event term is expanded as a sum of
the probabilities of the \(n\)-atoms that imply that event.  Nonnegativity,
total mass one, and the original constraints then form one finite linear
system.  A solution assigns mass to the corresponding atoms and hence
constructs a probability model.
The source's support lemma replaces a nonnegative solution of \(r\) linear
constraints by one having at most \(r\) positive entries.  Therefore
Theorem~2.4 needs at most \(|F|\) states for a satisfiable formula \(F\), and
Theorem~2.6 chooses rational masses of polynomial bit length.  Guessing that
finite model proves the upper half of Theorem~2.9; the reduction
\(A\mapsto w(A)>0\) from propositional satisfiability proves its lower half.
Richer conditional or nonlinear languages are not covered by that
\(\mathsf{NP}\)-completeness theorem.

\subsubsection{Baier--Katoen's finite DTMC, MDP, and PCTL presentation}

Chapter 10 of Baier and Katoen develops discrete-time Markov chains, cylinder
measures, reachability, PCTL and PCTL*, cost extensions, Markov decision
processes, schedulers, and optimal reachability.  The finite-state PCTL
algorithms label states bottom up: bounded until uses a finite recurrence;
unbounded until first separates value-zero and value-one regions and then
solves linear or Bellman equations.  These algorithms require a finite model;
the existence of a path measure on a general chain is not itself a finite
algorithm \cite[ch.~10]{BaierKatoen2008}.

For example, let \(g\) and \(d\) be absorbing states.  From \(u\), move to \(g\)
with probability \(1/3\) and to \(d\) with probability \(2/3\); from \(s\),
move to \(g\) or \(u\), each with probability \(1/2\).  With \(g\) as the target,
the graph just described has, in the state order \((s,u,g,d)\), transition
matrix
\[
 P=
 \begin{pmatrix}
  0&1/2&1/2&0\\
  0&0&1/3&2/3\\
  0&0&1&0\\
  0&0&0&1
 \end{pmatrix}.
\]
For instance,
\(\Pr_s(\operatorname{Cyl}(sug))=(1/2)(1/3)=1/6\).  The direct
one-step cylinder \(sg\) has mass \(1/2\).  Accordingly,
the unbounded reachability equations give
\[
 x_g=1,\qquad x_d=0,\qquad x_u=\frac13,\qquad
 x_s=\frac12+\frac12x_u=\frac23.
\]
The probability of reaching \(g\) within one transition from \(s\) is only
\(1/2\); the unbounded value \(2/3\) includes the two-step cylinder through
\(u\).  Thus \(s\models\mathsf P_{\geq2/3}[\top\mathsf U g]\), while
\(s\not\models\mathsf P_{\geq2/3}[\top\mathsf U^{\leq1}g]\).

To expose the scheduler quantifier, make \(s\) an MDP state with two actions.
Action \(a\) reaches \(g\) with probability \(3/4\) and \(d\) with probability
\(1/4\).  Action \(b\) reaches \(g\) with probability \(1/4\) and returns to
\(s\) with probability \(3/4\).  Always taking \(a\) gives reachability
\(3/4\); always taking \(b\) gives probability \(1\), because the probability
of returning forever is \(\lim_{n\to\infty}(3/4)^n=0\).  The Bellman equations
therefore yield
\[
 \Pr_s^{\min}(\Diamond g)=\frac34,
 \qquad
 \Pr_s^{\max}(\Diamond g)=1.
\]
The same graph consequently satisfies different minimum- and maximum-PCTL
threshold formulas.  This is why an undecorated probability operator cannot be
used on an MDP.

PCTL* enlarges the path language by allowing Boolean and temporal nesting and
is handled in the source by automata-theoretic product constructions rather
than by merely reusing the PCTL recurrence.  The source also develops chains
and decision processes with accumulated costs.  Those are separate extensions:
neither the PCTL* product construction nor a cost-bounded value algorithm is
claimed by Theorem~\ref{thm:p09-pctl-model-checking}.

The registered book contains no CTMC or CSL section.  Its DTMC/MDP/PCTL
presentation may not be cited for continuous-time claims.

\subsubsection{The separate CTMC and CSL authority}

Baier, Haverkort, Hermanns, and Katoen use finite rate-labelled CTMCs.  Their
Definition~1 types the rate map on all of \(S\times S\), permits positive rate
self-loops, and puts \(E(s)=\sum_tR(s,t)\).  Definitions~3--6 construct timed
paths and their measure and then give the CSL state/path grammar and semantics.
Proposition~2 handles the steady-state operator by bottom strongly connected
components and linear equations; Corollary~1 gives the embedded-chain equations
for unbounded until; and Theorem~1 gives the least-fixed-point Volterra system
for time-bounded until.  Theorems~2--3 in \S4.2 reduce the bounded intervals
\([0,t]\) and \([a,b]\) to one or two transient analyses on explicitly
transformed CTMCs \cite{BaierHaverkortHermannsKatoen2003}.

In the one-transition chain with rate \(R(s,g)=2\) and absorbing target \(g\),
the residence time in \(s\) has distribution
\[
 \Pr_s(T\leq t)=1-e^{-2t}.
\]
Hence the probability of satisfying
\(\top\mathsf U^{[0,1]}g\) is \(1-e^{-2}\), approximately \(0.8647\).  A
calculation may conclude the threshold \(3/4\) only from an approximation
interval whose lower endpoint exceeds \(3/4\); the decimal alone is not a
proof.

Section~4.3 supplies the exact approximation locator.  With generator
\(Q=R-\operatorname{diag}(E)\), uniformization rate
\(0<\lambda\) satisfying \(\lambda\geq\max_sE(s)\), and stochastic matrix
\(U=I+Q/\lambda\), equation~(7) expands the transient vector as a Poisson-
weighted sum of \(\mu U^n\).  Equation~(8) chooses the truncation index so that
the omitted Poisson mass is at most the requested \(\varepsilon\).  This bounds
truncation error; it does not by itself bound roundoff in the retained terms or
decide equality with a threshold.  Section~6 then gives the operator-by-
operator cost analysis.  Those are the claims used in
Theorem~\ref{thm:p09-csl-model-checking}; no CTMC claim is imported from the
discrete-time book.
